\ifdefined\gkver\else\def\gkver{student}\fi
\documentclass[\gkver]{gaokaozhenti}
\begin{document}

\chapter{2007年山东理综}
%% paperType: 理综物理部分

\begin{enumerate}
\item
%% number: 16
%% paperName: 2007年山东理综第 16 题 4 分
%% typeId: 1
%% type: 单选题
%% chapter: 力物体的平衡
%% point: 受力分析
%% method: 无
%% score: 4
%% degree: 650
%% duplicateId: 0
%% body:
如图所示，物体 A 靠在竖直墙面上，在力 $F$ 作用下，A、B 保持静止。物体 B 的受力个数为\xzanswer{C}

\onepicture[w=2.6cm]{figs/16.png}

\fourchoices[answer=C]
{2}
{3}
{4}
{5}

%% answer: C
%% memo:
\memoanswer{
本题原卷及材料中未提供详解，待补充。
}

\item
%% number: 17
%% paperName: 2007年山东理综第 17 题 4 分
%% typeId: 2
%% type: 多选题
%% chapter: 运动和力
%% point: 超重失重
%% method: 无
%% score: 4
%% degree: 650
%% duplicateId: 0
%% body:
下列实例属于超重现象的是\xzanswer{BD}

\fourchoices[answer=BD]
{汽车驶过拱形桥顶端}
{荡秋千的小孩通过最低点}
{跳水运动员被跳板弹起，离开跳板向上运动}
{火箭点火后加速升空}

%% answer: BD
%% memo:
\memoanswer{
本题原卷及材料中未提供详解，待补充。
}

\item
%% number: 18
%% paperName: 2007年山东理综第 18 题 4 分
%% typeId: 1
%% type: 单选题
%% chapter: 交流电
%% point: 变压器
%% method: 无
%% score: 4
%% degree: 650
%% duplicateId: 0
%% body:
某变压器原、副线圈匝数比为 55:9，原线圈所接电源电压按图示规律变化，副线圈接有负载。下列判断正确的是\xzanswer{D}

\onepicture[w=7.5cm]{figs/18.png}

\fourchoices[answer=D]
{输出电压的最大值为 $36\UV$}
{原、副线圈中电流之比为 55:9}
{变压器输入、输出功率之比为 55:9}
{交流电源有效值为 $220\UV$，频率为 $50\UHz$}

%% answer: D
%% memo:
\memoanswer{
本题原卷及材料中未提供详解，待补充。
}

\item
%% number: 19
%% paperName: 2007年山东理综第 19 题 4 分
%% typeId: 2
%% type: 多选题
%% chapter: 电场
%% point: 电场线
%% method: 无
%% score: 4
%% degree: 650
%% duplicateId: 0
%% body:
如图所示，某区域电场线左右对称分布，M、N 为对称线上两点。下列说法正确的是\xzanswer{AC}

\onepicture[w=6.0cm]{figs/19.png}

\fourchoices[answer=AC]
{M 点电势一定高于 N 点电势}
{M 点场强一定大于 N 点场强}
{正电荷在 M 点的电势能大于在 N 点的电势能}
{将电子从 M 点移动到 N 点，电场力做正功}

%% answer: AC
%% memo:
\memoanswer{
本题原卷及材料中未提供详解，待补充。
}

\item
%% number: 20
%% paperName: 2007年山东理综第 20 题 4 分
%% typeId: 1
%% type: 单选题
%% chapter: 运动和力
%% point: 牛顿定律中的图象
%% method: 无
%% score: 4
%% degree: 450
%% duplicateId: 0
%% body:
如图所示，光滑轨道 MO 和 ON 底端对接且 $ON=2MO$，M、N 两点高度相同。小球自 M 点由静止自由滚下，忽略小球经过 O 点时的机械能损失，以 $v$、$s$、$a$、$E_{k}$ 分别表示小球的速度、位移、加速度和动能四个物理量的大小。下列图象中能正确反映小球自 M 点到 N 点运动过程的是\xzanswer{A}

\onepicture[w=4.5cm]{figs/20.png}

\fourchoices[answer=A, ispicture=true, h=2.2cm]
{figs/20a.png}
{figs/20b.png}
{figs/20c.png}
{figs/20d.png}

%% answer: A
%% memo:
\memoanswer{
本题原卷及材料中未提供详解，待补充。
}

\item
%% number: 21
%% paperName: 2007年山东理综第 21 题 4 分
%% typeId: 1
%% type: 单选题
%% chapter: 电磁感应
%% point: 电磁感应电路分析
%% method: 无
%% score: 4
%% degree: 450
%% duplicateId: 0
%% body:
用相同导线绕制的边长为 $L$ 或 $2L$ 的四个闭合导体线框，以相同的速度匀速进入右侧匀强磁场，如图所示。在每个线框进入磁场的过程中，M、N 两点间的电压分别为 $U_{a}$、$U_{b}$、$U_{c}$ 和 $U_{d}$。下列判断正确的是\xzanswer{B}

\onepicture[w=11.0cm]{figs/21.png}

\fourchoices[answer=B]
{$U_{a}<U_{b}<U_{c}<U_{d}$}
{$U_{a}<U_{b}<U_{d}<U_{c}$}
{$U_{a}=U_{b}<U_{c}=U_{d}$}
{$U_{b}<U_{a}<U_{d}<U_{c}$}

%% answer: B
%% memo:
\memoanswer{
本题原卷及材料中未提供详解，待补充。
}

\item
%% number: 22
%% paperName: 2007年山东理综第 22 题 4 分
%% typeId: 2
%% type: 多选题
%% chapter: 曲线运动
%% point: 人造地球卫星
%% method: 无
%% score: 4
%% degree: 450
%% duplicateId: 0
%% body:
2007 年 4 月 24 日，欧洲科学家宣布在太阳之外发现了一颗可能适合人类居住的类地行星 Gliese581c。这颗围绕红矮星 Gliese581 运行的星球有类似地球的温度，表面可能有液态水存在，距离地球约为 20 光年，直径约为地球的 1.5 倍，质量约为地球的 5 倍，绕红矮星 Gliese581 运行的周期约为 13 天。假设有一艘宇宙飞船飞临该星球表面附近轨道，下列说法正确是\xzanswer{BC}

\fourchoices[answer=BC]
{飞船在 Gliese581c 表面附近运行的周期约为 13 天}
{飞船在 Gliese581c 表面附近运行时的速度大于 $7.9\Ukms$}
{人在 Gliese581c 上所受重力比在地球上所受重力大}
{Gliese581c 的平均密度比地球平均密度小}

%% answer: BC
%% memo:
\memoanswer{
本题原卷及材料中未提供详解，待补充。
}

\item
%% number: 23
%% paperName: 2007年山东理综第 23 题 11 分
%% typeId: 4
%% type: 实验题
%% chapter: 电路
%% point: 测量电阻率
%% method: 无
%% score: 11
%% degree: 450
%% duplicateId: 0
%% body:
检测一个标称值为 $5\UO$ 的滑动变阻器。可供使用的器材如下：

A．待测滑动变阻器 $R_{x}$，全电阻约 $5\UO$（电阻丝绕制紧密，匝数清晰可数）

B．电流表 $\mathrm{A}_{1}$，量程 $0.6\UA$，内阻约 $0.6\UO$

C．电流表 $\mathrm{A}_{2}$，量程 $3\UA$，内阻约 $0.12\UO$

D．电压表 $\mathrm{V}_{1}$，量程 $15\UV$，内阻约 $15\UkO$

E．电压表 $\mathrm{V}_{2}$，量程 $3\UV$，内阻约 $3\UkO$

F．滑动变阻器 $R$，全电阻约 $20\UO$

G．直流电源 $E$，电动势 $3\UV$，内阻不计

H．游标卡尺

I．毫米刻度尺

J．电键 S、导线若干

\begin{enumerate}
	\item
	用伏安法测定 $R_{x}$ 的全电阻值，所选电流表 \tkanswer[1.2cm]{$\mathrm{A}_{1}$}（填“A$_{1}$”或“A$_{2}$”），所选电压表为 \tkanswer[1.2cm]{$\mathrm{V}_{2}$}（填“V$_{1}$”或“V$_{2}$”）。

	\item
	画出测量电路的原理图，并根据所画原理图将下图中实物连接成测量电路。

	\onepicture[w=6.0cm]{figs/23.png}

	\item
	为了进一步测量待测量滑动变阻器电阻丝的电阻率，需要测量电阻丝的直径和总长度，在不破坏变阻器的前提下，请设计一个实验方案，写出所需器材及操作步骤，并给出直径和总长度的表达式。
\end{enumerate}

%% answer:
\jdanswer{
\begin{enumerate}
	\item
	$\mathrm{A}_{1}$；$\mathrm{V}_{2}$
	\item
	测量电路的原理图和实物连接图如图所示（方案一：分压接法；方案二：限流接法）
	\item
	见详解
\end{enumerate}
}

%% memo:
\memoanswer{
（1）$\mathrm{A}_{1}$；$\mathrm{V}_{2}$。

（2）电路原理图和对应的实物连接如图所示。

方案一：分压接法

\onepicture[w=4.0cm]{figs/23_an1.png}

\onepicture[w=6.0cm]{figs/23_an3.png}

方案二：限流接法

\onepicture[w=4.0cm]{figs/23_an2.png}

\onepicture[w=6.0cm]{figs/23_an4.png}

（3）方案一：

需要的器材：游标卡尺、毫米刻度尺

主要操作步骤：

①数出变阻器线圈缠绕匝数 $n$；

②用毫米刻度尺（也可以用游标卡尺）测量所有线圈的排列长度 $L$，可得电阻丝的直径为 $d=\frac{L}{n}$；

③用游标卡尺测量变阻器线圈部分的外径 $D$，可得电阻丝总长度 $l=n\pi\left(D-\frac{L}{n}\right)$；也可以用游标卡尺测量变阻器瓷管部分的外径 $D$，得电阻丝总长度 $l=n\left(D-\frac{L}{n}\right)$；

④重复测量三次，求出电阻丝直径和总长度的平均值。

方案二：

需要的器材：游标卡尺

主要的操作步骤：

①数出变阻器线圈缠绕匝数 $n$；

②用游标卡尺测量变阻器线圈部分的外径 $D_{1}$ 和瓷管部分的外径 $D_{2}$，可得电阻丝的直径为 $d=\frac{D_{1}-D_{2}}{2}$，电阻丝总长度 $l=\frac{n}{2}\pi(D_{1}+D_{2})$；

③重复测量三次，求出电阻丝直径和总长度的平均值。
}

\item
%% number: 24
%% paperName: 2007年山东理综第 24 题 16 分
%% typeId: 6
%% type: 计算题
%% chapter: 曲线运动
%% point: 圆周运动的应用
%% method: 无
%% score: 16
%% degree: 450
%% duplicateId: 0
%% body:
如图所示，一水平圆盘绕过圆心的竖直轴转动，圆盘边缘有一质量 $m=1.0\Ukg$ 的小滑块。当圆盘转动的角速度达到某一数值时，滑块从圆盘边缘滑落，经光滑的过渡圆管进入轨道 ABC。已知 AB 段斜面倾角为 $53^{\circ}$，BC 段斜面倾角为 $37^{\circ}$，滑块与圆盘及斜面间的动摩擦因数均 $\mu=0.5$，A 点离 B 点所在水平面的高度 $h=1.2\Um$。滑块在运动过程中始终未脱离轨道，不计在过渡圆管处和 B 点的机械能损失，最大静摩擦力近似等于滑动摩擦力，取 $g=10\Umsq$，$\sin37^{\circ}=0.6$，$\cos37^{\circ}=0.8$。

\begin{enumerate}
	\item
	若圆盘半径 $R=0.2\Um$，当圆盘的角速度多大时，滑块从圆盘上滑落？
	\item
	若取圆盘所在平面为零势能面，求滑块到达 B 点时的机械能。
	\item
	从滑块到达 B 点时起，经 $0.6\Us$ 正好通过 C 点，求 BC 之间的距离。
\end{enumerate}

\onepicture[w=8.0cm, align=right]{figs/24.png}

%% answer:
\jdanswer{
\begin{enumerate}
	\item
	$5\Urads$
	\item
	$-4\UJ$
	\item
	$0.76\Um$
\end{enumerate}
}

%% memo:
\memoanswer{
（1）滑块在圆盘上做圆周运动时，静摩擦力充当向心力，根据牛顿第二定律，可得：
\[ \mu mg=m\omega^{2}R \]
代入数据解得：$\omega=\sqrt{\frac{\mu g}{R}}=5\Urads$。

（2）滑块在 A 点时的速度：$v_{A}=\omega R=1\Ums$。

从 A 到 B 的运动过程由动能定理得：
\[ mgh-\mu mg\cos53^{\circ}\times\frac{h}{\sin53^{\circ}}=\frac{1}{2}mv_{B}^{2}-\frac{1}{2}mv_{A}^{2} \]
在 B 点时的机械能为：$E_{B}=\frac{1}{2}mv_{B}^{2}-mgh=-4\UJ$。

（3）滑块在 B 点时的速度：$v_{B}=4\Ums$。

滑块沿 BC 段向上运动时的加速度大小：$a_{1}=g(\sin37^{\circ}+\mu\cos37^{\circ})=10\Umsq$。

返回时的加速度大小：$a_{2}=g(\sin37^{\circ}-\mu\cos37^{\circ})=2\Umsq$。

BC 间的距离：$s_{BC}=\frac{v_{B}^{2}}{2a_{1}}-\frac{1}{2}a_{2}\left(t-\frac{v_{B}}{a_{1}}\right)^{2}=0.76\Um$。
}

\item
%% number: 25
%% paperName: 2007年山东理综第 25 题 18 分
%% typeId: 6
%% type: 计算题
%% chapter: 磁场
%% point: 带电粒子在复合场中的运动
%% method: 无
%% score: 18
%% degree: 450
%% duplicateId: 0
%% body:
飞行时间质谱仪可以对气体分子进行分析。如图所示，在真空状态下，脉冲阀 P 喷出微量气体，经激光照射产生不同价位的正离子，自 a 板小孔进入 a、b 间的加速电场，从 b 板小孔射出，沿中线方向进入 M、N 板间的偏转控制区，到达探测器。已知元电荷电量为 $e$，a、b 板间距为 $d$，极板 M、N 的长度和间距均为 $L$。不计离子重力及进入 a 板时的初速度。

\begin{enumerate}
	\item
	当 a、b 间的电压为 $U_{1}$ 时，在 M、N 间加上适当的电压 $U_{2}$，使离子到达探测器。请导出离子的全部飞行时间与比荷 $K$（$K=ne/m$）的关系式。
	\item
	去掉偏转电压 $U_{2}$，在 M、N 间区域加上垂直于纸面的匀强磁场，磁感应强度 $B$，若进入 a、b 间所有离子质量均为 $m$，要使所有的离子均能通过控制区从右侧飞出，a、b 间的加速电压 $U_{1}$ 至少为多少？
\end{enumerate}

\onepicture[w=8.5cm, align=right]{figs/25.png}

%% answer:
\jdanswer{
\begin{enumerate}
	\item
	$t=\frac{2d+L}{\sqrt{2KU_{1}}}$
	\item
	$U_{\min}=\frac{25eL^{2}B^{2}}{32m}$
\end{enumerate}
}

%% memo:
\memoanswer{
（1）由动能定理：
\[ neU_{1}=\frac{1}{2}mv^{2} \]
n 价正离子在 a、b 间的加速度：
\[ a_{1}=\frac{neU_{1}}{md} \]
在 a、b 间运动的时间：
\[ t_{1}=\frac{v}{a_{1}}=\sqrt{\frac{2m}{neU_{1}}}d \]
在 MN 间运动的时间：
\[ t_{2}=\frac{L}{v} \]
离子到达探测器的时间：$t=t_{1}+t_{2}=\frac{2d+L}{\sqrt{2KU_{1}}}$。

（2）假定 n 价正离子在磁场中向 N 板偏转，洛伦兹力充当向心力，设轨迹半径为 $R$，由牛顿第二定律得：
\[ nevB=m\frac{v^{2}}{R} \]
离子刚好从 N 板右侧边缘穿出时，由几何关系：
\[ R^{2}=L^{2}+(R-L/2)^{2} \]
由以上各式得：
\[ U_{1}=\frac{25neL^{2}B^{2}}{32m} \]
当 $n=1$ 时 $U_{1}$ 取最小值 $U_{\min}=\frac{25eL^{2}B^{2}}{32m}$。
}

\item
%% number: 36
%% paperName: 2007年山东理综第 36 题 8 分
%% typeId: 6
%% type: 计算题
%% chapter: 气体性质
%% point: 查理定律
%% method: 无
%% score: 8
%% degree: 450
%% duplicateId: 0
%% body:
某压力锅结构如图所示。盖好密封锅盖，将压力阀套在出气孔上，给压力锅加热，当锅内气体压强达到一定值时，气体就把压力阀顶起。假定在压力阀被顶起时，停止加热。

\begin{enumerate}
	\item
	若此时锅内气体的体积为 $V$，摩尔体积为 $V_{0}$，阿伏加德罗常数为 $N_{A}$，写出锅内气体分子数的估算表达式。
	\item
	假定在一次放气过程中，锅内气体对压力阀及外界做功 $1\UJ$，并向外界释放了 $2\UJ$ 的热量。锅内原有气体的内能如何变化？变化了多少？
	\item
	已知大气压强 $p$ 随海拔高度 $H$ 的变化满足 $p=p_{0}(1-\alpha H)$，其中常数 $\alpha>0$。结合气体定律定性分析在不同的海拔高度使用压力锅，当压力阀被顶起时锅内气体的温度有何不同。
\end{enumerate}

\onepicture[w=5.5cm, align=right]{figs/36.png}

%% answer:
\jdanswer{
\begin{enumerate}
	\item
	$\frac{V}{V_{0}}N_{A}$
	\item
	内能减少了 $3\UJ$
	\item
	阀门被顶起时锅内气体温度随着海拔高度的增加而降低
\end{enumerate}
}

%% memo:
\memoanswer{
（1）设锅内气体分子数为 $n$，则 $n=\frac{V}{V_{0}}N_{A}$。

（2）根据热力学第一定律得：$\Delta U=W+Q=-3\UJ$。锅内气体内能减少，减少了 $3\UJ$。

（3）由 $p=p_{0}(1-\alpha H)$（其中 $\alpha>0$）知，随着海拔高度的增加，大气压强减小。由 $p_{1}=p_{0}+\frac{mg}{S}$ 知，随着海拔高度的增加，阀门被顶起时锅内气体压强减小。根据查理定律得 $\frac{p_{1}}{T_{1}}=\frac{p_{2}}{T_{2}}$，可知阀门被顶起时锅内气体温度随着海拔高度的增加而降低。
}

\item
%% number: 37
%% paperName: 2007年山东理综第 37 题 8 分
%% typeId: 6
%% type: 计算题
%% chapter: 机械波
%% point: 波长频率波速
%% method: 无
%% score: 8
%% degree: 650
%% duplicateId: 0
%% body:
湖面上一点 O 上下振动，振幅为 $0.2\Um$，以 O 点为圆心形成圆形水波，如图所示，A、B、O 三点在一条直线上，OA 间距离为 $4.0\Um$，OB 间距离为 $2.4\Um$。某时刻 O 点处在波峰位置，观察发现 $2\Us$ 后此波峰传到 A 点，此时 O 点正通过平衡位置向下运动，OA 间还有一个波峰。将水波近似为简谐波。

\begin{enumerate}
	\item
	求此水波的传播速度、周期和波长。
	\item
	以 O 点处在波峰位置为 0 时刻，某同学打算根据 OB 间距离与波长的关系，确定 B 点在 0 时刻的振动情况，画出 B 点的振动图像。你认为该同学的思路是否可行？若可行，画出 B 点振动图像，若不可行，请给出正确思路并画出 B 点的振动图象。
\end{enumerate}

\onepicture[w=3.5cm, align=right]{figs/37.png}

%% answer:
\jdanswer{
\begin{enumerate}
	\item
	$v=2\Ums$，$T=1.6\Us$，$\lambda=3.2\Um$
	\item
	可行，B 点的振动图象如图
\end{enumerate}
}

%% memo:
\memoanswer{
（1）$v=\frac{\Delta x_{1}}{\Delta t}=2\Ums$；$\Delta t=\frac{5}{4}T$，解得 $T=\frac{4}{5}\Delta t=1.6\Us$；$\lambda=vT=3.2\Um$。

（2）可行。振动图象如图。

\onepicture[w=6.0cm]{figs/37_an.png}
}

\item
%% number: 38
%% paperName: 2007年山东理综第 38 题 4 分
%% typeId: 6
%% type: 计算题
%% chapter: 运动和力
%% point: 动量和能量
%% method: 无
%% score: 4
%% degree: 650
%% duplicateId: 0
%% body:
人类认识原子结构和开发利用原子能经历了十分曲折的过程。请按要求回答下列问题。

\begin{enumerate}
	\item
	卢瑟福、玻尔、查德威克等科学家在原子结构或原子核的研究方面做出了卓越的贡献。请选择其中的两位，指出他们的主要成绩。

	①\tkanswer[4.0cm]{卢瑟福提出了原子的核式结构模型（或其他成就）}；

	②\tkanswer[4.0cm]{玻尔把量子理论引入原子模型，并成功解释了氢光谱（或其他成就）}。

	在贝克勒尔发现天然放射现象后，人们对放射线的性质进行了深入研究，下图为三种射线在同一磁场中的运动轨迹，请从三种射线中任选一种，写出它的名称和一种用途。

	\tkanswer[6.0cm]{①高速 $\alpha$ 粒子流可以用在人工核转变中；②$\gamma$ 光子可以用于医疗放射治肿瘤}。

	\onepicture[w=3.2cm, align=right]{figs/38.png}

	\item
	在可控核反应堆中需要给快中子减速，轻水、重水和石墨等常用作减速剂。中子在重水中可与 $\ce{^{2}_{1}H}$ 核碰撞减速，在石墨中与 $\ce{^{12}_{6}C}$ 核碰撞减速。上述碰撞可简化为弹性碰撞模型。某反应堆中快中子与静止的靶核发生对心正碰，通过计算说明，仅从一次碰撞考虑，用重水和石墨作减速剂，哪种减速效果更好？
\end{enumerate}

%% answer:
\jdanswer{
\begin{enumerate}
	\item
	卢瑟福提出了原子的核式结构模型（或其他成就）；玻尔把量子理论引入原子模型，并成功解释了氢光谱（或其他成就）；高速 $\alpha$ 粒子流可以用在人工核转变中，$\gamma$ 光子可以用于医疗放射治肿瘤。
	\item
	与重水靶核碰后中子速度较小，故重水减速效果更好。
\end{enumerate}
}

%% memo:
\memoanswer{
设中子质量为 $M_{n}$，靶核质量为 $M$。

由动量守恒定律：$M_{n}v_{0}=M_{n}v_{1}+Mv_{2}$。

由能量守恒定律：$\frac{1}{2}M_{n}v_{0}^{2}=\frac{1}{2}M_{n}v_{1}^{2}+\frac{1}{2}Mv_{2}^{2}$。

解得：
\[ v_{1}=\frac{M_{n}-M}{M_{n}+M}v_{0} \]
在重水中靶核质量：$M_{H}=2M_{n}$，
\[ v_{1H}=\frac{M_{n}-M_{H}}{M_{n}+M_{H}}v_{0}=-\frac{1}{3}v_{0} \]
在石墨中靶核质量：$M_{C}=12M_{n}$，
\[ v_{1C}=\frac{M_{n}-M_{C}}{M_{n}+M_{C}}v_{0}=-\frac{11}{3}v_{0} \]
与重水靶核碰后中子速度较小，故重水减速效果更好。
}

\end{enumerate}

\end{document}
